\subsection{The adelic division pair on the completed point}
\label{subsec:completed-division}

The preceding period comparison identifies the operator in the ramified coordinates. For the division argument this identification must also respect the degree-zero forms, their metrics and the refinements of level at every word.

Fix a completed member \(q\) of Definition~\ref{def:completed-member}
and its compatible reciprocity data \(R\). For each word \(u\) the
degree-zero column forms of \(R\) agree with those of \(q\). Apply this equality to the measured
saturation and its remaining spectral data. This gives the
measured saturation for \(R\) with the same group, measure,
evaluation translate and comparison of the full and good rows.
The finite part of the saturation puts each period indexed by the sequence in the
level space of its indicated character. Thus the period
construction above uses the prescribed level projector on
every word.

\begin{proposition}[The completed division pair]\label{prop:completed-division-pair}
On a completed member in the sense of Definition~\ref{def:completed-member}, the real, Hecke, and lower
comparisons identify the square-integrable and geometric
realizations of the definite columns, periods, and size
functions at all subspaces. The identified point carries
the original tuple characters and reciprocity data,
including the prescribed side at each selected prime.
In particular, writing \(q,R\) for the two identified realizations and \(F_u^0\) for the degree-zero forms of word \(u\),
\[
 F_u^0(q)=F_u^0(R)\qquad\text{for every }u.
\]
Both parameter runs are evaluated on this same completed point.
\end{proposition}

\begin{proof}
Use the two ambient identifications on \(q\), the intermediate
modules and the finite group comparison of the preceding
construction. Denote the resulting adelic row data and face by
\(x\). The adelic division pair is formed from the base and
sharp form readings of \(x\) and its word columns. Its
intersection comparison over the coefficient field is the
base change of the one on the original row. The rational word
isomorphisms and the sharp and base ambient identifications
are those in the construction of \(x\).

Let \(n_{\rm sh}\) be the sharp-level index. Write \(E\) for
sharp evaluation with the measured translate and Petersson
comparison, using the conjugate scalar action at a conjugate
reading. The placement chosen in the preceding construction
is \(s_{\rm sh}=\kappa E\) with \(\kappa>0\) and
\[
                 n_{\rm sh}\|s_{\rm sh}v\|^2=\|v\|^2
\]
on the refined row. In particular its phase is \(1\). Writing \(F_u^0\) for the degree-zero form space, the identification on the two realizations is wordwise:
\[
 F_u^0(q)=F_u^0(R)\qquad\text{for each word }u .
\]
The saturation therefore uses the same level projector, character,
evaluation translate and Haar measure in the norm identity.

Let \(e_F,\mathrm{fi}\) be the two coefficient maps in the normal
expansion. A substitution of face sizes or error changes the
arguments of the expansion. Its form and polynomial coefficients
are read in the original rational frame, so the coefficient maps
are still \(e_F,\mathrm{fi}\). The diagonal forms, positive
terms, standard Hermitian frames and their lattices are those
of \(q\); uniqueness gives the standard additive character of
sign \(-1\) on both sides. Use the bounded real comparison on
this core with its frame weights, extremal vectors, level action
and local uniformization maps. All its phase factors are the
one just computed. This identifies the intersection, norm and
period comparisons on a single pair. \label{cor:completed-division-1}

\proofstep{Identify the actual Hecke eigenvalues and tuple characters.}

On the base row write the decomposition by complete Hecke systems as
\[
 V_{\rm base}=\bigoplus_{\lambda}V_{\rm base}[\lambda].
\] Fix one
occurring eigensystem \(\lambda\) and a required tuple \(\mu\).
Multiplicity gives its nonzero constituent at \(\mu\) with
the prescribed infinite components off the distinguished place.
The generic series choice makes the product of the finite
character and level projectors nonzero on the required labels.
At each tuple in use the adelic and marked characters satisfy
\[
 \chi^{\rm adelic}_\mu=\chi^{\rm marked}_\mu=\chi^R_\mu .
\]
The first equality uses the tuple map and the second uses
compatibility with \(R\). In particular the character depends only on the tuple label. For two positions \(a,b\) the implication is
\[
 \mu(a)=\mu(b)\quad\Longrightarrow\quad \chi_{\mu(a)}=\chi_{\mu(b)} .
\] Discrete-series reciprocity at the actual
weights gives the relative parameter as the logarithm of the
positive lift parameter; intersection at the refined level gives
both rational realizations of this constituent. \label{cor:completed-division-2}

\proofstep{Identify the lower maps in both ambient realizations.}

At the base tuple use the reference lower map of the
standard columns. At a nonreference tuple the lower map uses
the integral columns of the actual twist. To see the
polynomial comparison there, let \(P\) be the highest-weight
polynomial in the span for the prescribed top weight. Take
the standard original and twisted columns and a last-stage
frame whose determinant is a unit.
Suppose the evaluation of the specified twisted vector on the
given class agrees at every argument with the function of \(P\)
in that frame, with its prescribed coordinate shift. In the standard last-stage coordinates these evaluations recover the polynomial
\[
 P_{\rm recovered}=\rho_\lambda(A)P,
\]
where \(A\) is the matrix of the last-stage frame and \(\rho_\lambda\) its polynomial action. Applying the lower
map means taking its branching projection. The two true
weights interlace.

In those weights the restriction to the
last block commutes with the projection in the sense
\[
                 \operatorname{br}\big(\rho_{\lambda}(A)P\big)
                  =\operatorname{br}_{\rm tens}(P)
\]
with the lower map expressed in its standard coordinates.
Substitution in the top formula followed by extraction in the standard twisted coordinates recovers \(\rho_\lambda(A)P\). The block branching formula, with the upper and
lower weight shifts, then gives the displayed equality and
the determinant factor of the lower coordinate is \(1\).
This comparison uses the equality of the twist datum, the
standard twisted column maps, and the standard original
columns on \(q\). The unit condition restricts the frame in this statement.

The sharp ambient map is the map obtained from the sharp
identification at \(q\) and the two forms of the tuple
reading. The splitting of the lower modules is the splitting
of the adelic pair. If \(s_1,s_2\) are the stabilizer orders at the two levels and \(C\) the common stabilizer multiple, the original class and level conditions give
\[
 s_1\mid C,\quad s_2\mid C,\qquad \operatorname{lcm}(s_1,s_2)\mid C .
\] These identities supply
both lower maps as well as the two ambient realizations on
\(x\). \label{cor:completed-division-3}

\proofstep{Identify the two sets of marked primes under the prescribed field isomorphism.}

Let \(\theta\) be the auxiliary field isomorphism from the
standard kernel to the column model. Write \(Q_{\ker}(\mu)\)
for its finite set of quotient primes. The two prime comparisons
in the construction give
\[
 \theta_*Q_{\ker}(\mu)=Q(\mu),\qquad
 \theta_*(\operatorname{im}\mathfrak P_{\ker})=
                         \operatorname{im}\mathfrak P .
\]
Here \(\mathfrak P(v)\) is the selected prime of \(E\) over
\(v\) in \(E^+\). A prime contracting under \(\theta\) to
\(P\in Q_{\ker}(\mu)\) is necessarily \(\theta_*P\). The
level comparison at a used tuple consequently puts that very
prime in the adelic quotient-prime set. For
\(\ell\in\{p_d,p_d'\}\), if \(\ell\in\mathfrak P(v)\), then
\(v\mid\ell\). The marked side law at \(v\) identifies the
contraction of \(\mathfrak P(v)\) to the imaginary quadratic
field with the \(\ell\)-prime specified by \(R\).
\label{cor:completed-division-4}

\proofstep{Substitute the common columns before comparing the periods and sizes.}
Compare periods and heights on these identified columns.


For any subspace \(U\), identify its square-integrable column
with its geometric column by the definite comparison. The
resulting period spaces on this one domain agree, including the
Poincar\'e period in the updated component. In symbols, writing
\(C,\mathcal P,H\) for the column, its period space and its size,
the successive identifications give
\[
 C_{L^2}(U)=C_{\rm geom}(U),\qquad
 \mathcal P_{L^2}(U)=\mathcal P_{\rm geom}(U),\qquad
 H_{L^2}(U)=H_{\rm geom}(U).
\]
For the last equality both sides use the record of \(q\), the
nonnegative integer values assigned to \(\Delta,\eta_{\rm sz}/m\) and the same
replacement of \(l_{6m,j},l_{6u}\) at the given real coefficient.
Transport of the bounded real datum through these identifications
preserves the weights, word placement, side and lower maps, on
the point with record \(R\). \label{cor:completed-division-5}
\end{proof}

\begin{construction}
\begin{tabular}{@{}l@{\hspace{1.2em}}l@{}}
\toprule
\textit{Datum of the completed point} & \textit{Identification used in both runs} \\
\midrule
Word and tuple character & Same rational column and level projector \\
Measured sharp evaluation \(E\) & \(s_{\rm sh}=\kappa E,\ \kappa>0\) \\
Sharp-level index \(n_{\rm sh}\) & \(n_{\rm sh}\|s_{\rm sh}v\|^2=\|v\|^2\) \\
Base and sharp forms & One adelic division pair at the chosen translate \\
\bottomrule
\end{tabular}
\caption{The comparison data carried into both parameter runs by Proposition~\ref{prop:completed-division-pair}.}\label{fig:analytic-completed-pair}
\end{construction}

\subsection{Two parameter runs on one division pair}
\label{subsec:completed-prime-runs}

Choose the point of Proposition~\ref{prop:completed-division-pair}.
The depth estimates and the exceptional-prime estimates call for two choices of numerical parameters. Their comparison uses the same row, level and period identification at this point.

\paragraph{Row normalization and the numerical point.}
The numerical point used for the runs has the face, word
lattices, period spaces and scale of the bundle. To form it,
take a recursive realization with the base and refined level
characters, a normalization of the real rows, and the local
row comparison between that normalization and the recursive
realization. This is the choice of comparison denoted by a
summand of the bundle. Denote the row spaces before and after this comparison by superscripts \(1,2\). The same identifications apply simultaneously to both rows and the word columns:
\[
 V_{\rm full}^{(1)}=V_{\rm full}^{(2)},\qquad V_{\rm good}^{(1)}=V_{\rm good}^{(2)},\qquad M_u^{(1)}=M_u^{(2)}\quad(\text{each word }u).
\]

The row normalization uses the equality of the character
frames, the identification of the type primes with the
kernel's type places and the conjugate-level and character
identities of the finite translate. It keeps the standard
frames and the coefficient values on the standard word
vectors. For the recursive comparison use the standard
additive character on these levels and the local row
reading of the marked Galois level. The degree-three
normalization uses the equality of the zero section with
the one in the original scale datum.

For \(z=v\) or \(\bar v\), write \(I_j(z),J_j(z),Z_j(z)\) for the
base-row index, lower-map index and value of the zero-section test at lift \(j\);
nonzero means that the corresponding test detects the place. The twist
support has the following intrinsic description:
\[
 \mathcal S_{\rm tw}
 =\{v:\ (I_j(z),J_j(z),Z_j(z))\ne(0,0,0)
            \text{ for some }j\text{ and }z\in\{v,\bar v\}\}.
\]
All three tests precede the choice of the height and error constants.
Let \(S\) be the set selected by the three tests and \(\mu_0\) the base tuple. The tests identify the tuple map on conjugate places by
\[
 S=\bar S,\qquad v\in S\ \Longleftrightarrow\ \mu(v)\ne\mu_0,\qquad \mu(\bar v)=\mu(v).
\] In particular each
single-index place is twisted.
These are comparisons on \(x\). Substitution of the three
family equalities of Section~\ref{subsec:completed-division}
transports them to the bundle with \(R\) fixed.
\label{cor:completed-division-6}

\paragraph{The row at the chosen comparison.}
\label{par:linked-family-rows}
Fix a completed member \(q\) and its division pair from
Proposition~\ref{prop:completed-division-pair}. Choose a recursive
realization, a normalization of its real rows and the identification
between these realizations on the pair. Write \(P\) for the resulting point. The row,
periods, finite groups, coefficient field and numerical functions
at \(P\) are those of the pair. In particular its two level rows
are the ones in Section~\ref{subsec:completed-division}. The
following verifications hold for each compatible comparison on
the bundle. Thus they also hold if an additional condition, such
as the one on the evaluation grid, selects the comparison.

\begin{proposition}[The two rows at a compatible choice of comparison]\label{prop:linked-comparison-rows}
At every compatible comparison on the completed bundle, let \(q_0\) be projection to the sharp good row and \(P_{u,0}\) the reference pairing of word \(u\). The factorization is characterized by
\[
 F-F'\in\ker q_0\quad\Longrightarrow\quad P_{u,0}(\phi,F)=P_{u,0}(\phi,F')\quad\text{for all }\phi.
\] At either level the rational row and its integral
submodule have the corresponding divisor comparisons of
the original division pair, on all tuples used by the real
row. Any later choice of a comparison satisfying an extra
condition therefore still gives these rows and their pairings.
\end{proposition}

\begin{proof}
\proofstep{Use the sharp reference period to test the kernel of the good projection.}
Fix a word \(u\), and use its selected Poincar\'e operator from
Section~\ref{subsec:selected-completion-period}. Let \(R_P\) be the
rational row, \(R_P^\sharp\) its common sharp good row and
\(q_0:R_P\to R_P^\sharp\) the good projection. In the complex
realization write \(I_u\) for the column isometry and \(I_R\) for
the isometry on the good row at the refined level. Unfolding by
Lemma~\ref{lem:joint-period} gives, for the sharp tests of this
operator,
\[
 \iota\bigl(P_{u,0}(\phi,F)\bigr)
  =\mathcal P_{u,0}\bigl(I_u\phi, I_R(q_0 F)\bigr),
\]
where \(\mathcal P_{u,0}\) denotes the resulting complex theta
pairing.

The normalized factor is one at the reference tuple
of Section~\ref{sec:transport}. The reference tuple is obtained
from the full diagonal and the dimension of the ramified space
there, before applying either division estimate. If \(F-F'\in\ker q_0\),
the difference of the right sides vanishes for every \(\phi\).
The coefficient embedding is injective, so the vanishing difference implies
\[
 F-F'\in\ker q_0\quad\Longrightarrow\quad P_{u,0}(\phi,F)=P_{u,0}(\phi,F')\qquad(\text{every }\phi).
\] This defines the rational
reference pairing on the quotient without choosing an integral
representative.

\proofstep{Read the quotient pairing on the transported integral lattices.}

The rational pairing obtained this way is the one characterized by
the Gram equations in Proposition~\ref{prop:rational-period-factor}.
For each of the two levels, the rational row comparison identifies
the down-lift and the theta map in the same coordinates. To see what
is integral, let \(M_w\) be the localization of the glued lattice
at a place \(w\).

For the coordinate identification \(a\) coming
from the completed pair, its lattice and form are transported by
\[
 M_w\longmapsto aM_w,\qquad
 B_w(\phi,F)\longmapsto B_w(a^{-1}\phi,a^{-1}F).
\]
Here \(a\) denotes the respective coordinate map in each argument.
Thus the localized orders, lattice indices and dual indices are
computed on the identified modules. The reference form descends
because of the kernel calculation above. A nonreference tuple of
the real row occurs as the tuple of a place obtained by the inverse
place bijection of the host. Hence for each used tuple \(\mu\) one of two comparisons is available:
\[
 \mu=\mu_0\quad\text{or}\quad \mu=\mu(w)\ \text{for a place }w\text{ on the nonreference support}.
\]
The measured intersection comparison supplies the surjections
onto the good coordinates and the separation of the primitive types.
The equality of the bounded real realizations identifies the real
datum and its face in both calculations.
\label{cor:linked-family-rows-1}

On a nonreference summand, let \(c\) be an integral level coordinate
reading to an admitted source value and \(\pi_{\rm src}(c)\) its
projection to the lift source. Write \(x=\Theta\pi_{\rm src}(c)\)
in form coordinates and choose \(x=\sum_h a_h f_h\), with \(f_h\)
rational forms in the expansion model and \(a_h\) in the row field.
If \({\rm ev}_w\) denotes the local expansion map, the divisor theta
identity and the signed series comparison give the same value
\[
        {\rm ev}_w(x)=\sum_h\iota_w(a_h)\,{\rm ev}_w(f_h).
\]

To check independence of the expansion, subtract the two expressions for \(x\) in the extended space and use linearity:
\[
 \sum_h b_h f_h=0\quad\Longrightarrow\quad \sum_h b_h\,{\rm ev}_w(f_h)=0 .
\] The completed
division comparison transports \(\pi_{\rm src}\), the rational form
coordinates and \({\rm ev}_w\) through the update of periods and
sizes. Consequently the displayed equality can be checked either
before or after the update. In the chosen rational basis, evaluation on the row is the composite 
\[
 x\longmapsto(a_h)_h\longmapsto(\iota_w(a_h))_h\longmapsto\sum_h \iota_w(a_h)\,{\rm ev}_w(f_h).
\]
 Its restrictions to the integral expansion lattice and theta sublattice give the integral evaluation and the assigned series, respectively.
\label{cor:linked-family-rows-2}

Forget the grid in the arithmetic member of
Section~\ref{subsec:automorphic-realization}. The norm choice,
the signed series used for the exponent, the two type maps, the
base and local level groups, the kernel field and the selected
half primes remain those of this member. The changes in the
Petersson scalar, in the lift scalars and in the error have no
effect on the maps of finite ideles or integral host rings.

Fix a preliminary frame and its finite support data. If \(S\)
is an allowed support, put
\[
    M(S)=g_0g_1\!\!\prod_{\substack{q\text{ a residue cardinality}\\
                         q\text{ of a place in }S}}\!q ,
\]
with each distinct cardinality used once, as in the level of the
member. By the support bound the possible \(M(S)\) belong to a
finite list. If \(A_M\) generates the embedded ray field at level \(M\), one finite request suffices:
\[
 A=\bigcup_{M\text{ in the level list}} A_M,\qquad |A|\le\sum_M |A_M|<\infty .
\] The Artin comparison of
Section~\ref{subsec:automorphic-realization} identifies these ray
fields and the type and character values in the chosen embedding.
The two canonical frames on the same core give the same half-level
ideals as the divisor, including the half-error divisibility.

The integral field reading identifies the coefficient rings. After
passing to fractions it identifies the degree and absolute
discriminant used in the size bound with those of the recorded core.
The single-index support test uses the case with zero error and
height constants. These comparisons give the field, rank, series
and level requirements in the two runs. A coordinate isomorphism preserves each of the module equalities in the comparison. For submodules \(U,V\) in an identified row and its isomorphism \(b\),
\[
 U=V\quad\Longleftrightarrow\quad b(U)=b(V).
\]
\label{cor:linked-family-rows-3}

At infinity the intertwiner is defined on the prescribed lowest
compact type of the polynomial Weil representation. Regularity,
admissibility and finite multiplicity give the Harish--Chandra
modules. Apply Casselman--Wallach globalization to this intertwiner
and cuspidal rapid decay to its integrals, as in
Section~\ref{sec:transport}. Let \(T\) be the difference of two such intertwiners,
and let \(V_0\) be the span of the compact translates of the
distinguished lowest vectors. Equality on the test vectors gives
\(T(V_0)=0\).

At the Harish--Chandra module level,
\[
 V=U(\mathfrak g_{\mathbb C})V_0,\qquad
 T(Xv)=X T(v)=0
 \quad\bigl(X\in U(\mathfrak g_{\mathbb C}),\ v\in V_0\bigr).
\]
Thus \(T=0\) on the module and, by continuity, on its globalization. This explains the injective reading used for the geometric
action. The conjugate frame is read by the transported coefficient
map. When \(n_V\ge8fr+4\), the stable theta span and the form
isomorphism of Section~\ref{subsec:selected-cusp-reading}
apply to these weights. Since the normalized map and the extended
intertwiner agree on every vector of the common orthonormal basis,
they agree on the row. In particular evaluation and the two
linear readings assign the same form at either frame.

At a supported finite type, let \(e_{\rm good}\) be the central
idempotent of Section~\ref{subsec:automorphic-realization}. For a
vector \(f\) in its image and a finite type-group element \(g\),
\[
       e_{\rm good}(g f)=g(e_{\rm good}f)=g f .
\]
The chosen nonzero component therefore generates a module for
the supported split factors in the same constituent. Apply
Green's regular cuspidal character and the depth-zero type
correspondence to the factor actions. Normalized induction and
uniqueness of cuspidal support give the multiset of the local
parameter. The marked local Galois identification reads it at
the coefficient choice of \(P\). This applies to each required
label with the threshold on the family of members. The primitive
base lifting isomorphism is rational and bijective in its
phase-one coordinates and has the uniform positive rank \(D_0\).
The tuple map, level subgroups and base character therefore
refer to this same row on the two conjugate sides.
\label{cor:linked-family-rows-4}

\end{proof}



The two runs below use the rows of Proposition~\ref{prop:linked-comparison-rows}
for the selected compatible comparison.

\begin{lemma}[The two prime orientations on the common core]
For \(\ell=p_d,p_d'\) there is a full short-section extension of
the same core \(J\) with distinguished prime \(\ell\) at every
short weight. For a twisted place \(v\), the determinant character and short cover use the marked half prime of the assigned orientation, where
\[
 p(v)=\begin{cases}p_d',&v\mid p_d,\\ p_d,&v\nmid p_d.\end{cases}
\]
\end{lemma}
\begin{proof}
Write \(J\) for the common geometric core at \(P\). It
has its canonical frames, positive ranks and half-level
divisibility of the error data. The two field conditions and
the values of both distinguished characters hold in the
prescribed fields. The half-level ideal \(\mathfrak m_b\)
of the common core is the product of the marked type prime
ideals of the row. In fact the full level is squarefree
and factors into its two indicated halves. The divisor
comparisons and the equality of the core and row half
primes give the same prime factors on the indicated
side. Equality follows from the squarefree factorization.

Fix \(\ell\in\{p_d,p_d'\}\). For this orientation choose
the distinguished character \(\delta_\ell\) of the frame at
\(\ell\), and use the field condition at \(\ell\). Define
the own prime of a coefficient place \(v\) by
\[
  p(v)=
  \begin{cases}
    p_d' & \text{if the residue characteristic of \(v\) is \(p_d\)},\\
    p_d & \text{otherwise}.
  \end{cases}
\]
For twisted places with \(p(v)=\ell\), the marked side
comparison supplies one prime of the imaginary quadratic
field over \(\ell\) for all the quotient primes over
\(\ell\) on the tuples of those places. Extension of that
prime is contained in every one of these quotient primes,
by contraction. If the class of such places is empty, any
prime above \(\ell\) suffices for this orientation.

On the finite ideles the character construction gives
characters \(\xi_\ell,d_\ell\) with
\[
                        \xi_\ell(a/\bar a)=d_\ell(a).
\]
Here \(d_\ell\) is the prescribed finite-idelic restriction of
\(\delta_\ell\). A character of finite order on the ideles
is trivial on the archimedean subgroup of the CM field,
since its archimedean part is a product of copies of \(\C^\times\).
For a character \(\xi\) of order \(n\) and any archimedean idele
\(z\), take an archimedean \(w\) with \(w^n=z\). Then
\[
 \xi(z)=\xi(w^n)=\xi(w)^n=1 .
\] Thus
its restriction to the finite ideles inherits triviality on principal ideles,
triviality on real ideles, and the specified coefficient
values.

The finite character is trivial on integral units congruent to \(1\)
at the prescribed level. At a prime over \(\ell\), choose
a unit with a nontrivial character value, using the
distinguished character's prescribed tame ramification.
An idele equal to this unit at that prime and to \(1\)
elsewhere is congruent to \(1\) at the two type primes
because their rational primes differ from \(\ell\).
This gives the required nontrivial test there.

On finite ideles put \(h(b)=b/\bar b\). Finite idele Hilbert
90 gives \(\operatorname{im}h=\ker N_{E/E^+}\). Since
\(\ker h\) consists of real ideles and \(d_\ell\) is trivial
there, the assignment
\[
 \xi_\ell(h(b))=d_\ell(b)
\]
is well defined. Extend this character to the finite idele
group; its required congruence and finite-order properties
concern the norm-one subgroup. For a principal norm-one
element \(a\), explicit preimages are
\[
 b=1+a\ (a\ne-1),\qquad
 b\in E^\times,\ \bar b=-b\ (a=-1).
\]
Thus \(\xi_\ell(a)=1\) on principal norm-one elements.
The congruence form of Hilbert 90 transports the tame
congruences at the two type primes and at \(\ell\), giving
the stipulated ramification and coefficient values.
\label{cor:completed-division-7}

Extend the core to the upper geometric model, then take the
full short-section extension with the prime and the
character just chosen. These constructions use the laws
of \(J\), the canonical frames, positive ranks and the
host hypotheses at \(\ell\). The full extension has the
same core \(J\), and its short cover at each lift has
distinguished prime \(\ell\) at every short weight.
For a determinant root \(b\) of an adelic element \(g\), let \(b_f\)
be its finite part. The finite determinants and character values satisfy
\[
 \det(g)_f=b_f/\bar b_f,\qquad
 \xi_\ell(\det(g)_f)=d_\ell(b_f).
\] Rational determinant roots give
value \(1\). At a twisted \(v\) with \(p(v)=\ell\), the
same formula gives the specified determinant character and
the chosen short cover above the marked half prime. This
verifies the own-prime and character pinning at that place.
Taking \(\ell=p_d\) and \(\ell=p'_d\) gives the two extensions.
\label{cor:completed-division-8}
\end{proof}

\begin{established}[The split refinement of the marked core]\label{hyp:split-section-refinement}
In the local construction write \(p=\ell\), \(M=g_0g_1\) for
the original full level and \(I=P\mathcal O_E\), where \(P\)
is the chosen prime above \(p\) in the marked imaginary
quadratic field. If \(J\) denotes the old half, then the new level data are
\[
 (\text{level},\text{chosen half},\text{auxiliary level},\text{field})
                       =(Mp,JI,1,F),
\]
with the same coefficient field \(F\). For each \(\sigma\in\operatorname{Aut}(\mathbb C/F)\), the marked CM curve and its pullback are isomorphic by an isomorphism respecting every torsion label at level \(Mp\) and at each divisor of \(Mp\).
\end{established}
\begin{proof}
We have \(I+\overline I=\mathcal O_E\) and
\(I\overline I=(p)\), by splitting. Use at each lift the
polarization and moduli problem of the core \(J\) just constructed.
At each \(j\) refine the split level by \(I,\overline I\).
Writing \(J\) for the old chosen half, the updated parameters are
\[
 (\text{integer},\text{chosen half},\text{auxiliary level})=(Mp,JI,1).
\] The two complementary halves exist by the displayed
factorization and coprimality. Use the moduli space with this
refined marking and the map forgetting its new level.
The coefficient field stays fixed. The field contains the specified roots of the coprime orders \(M\) and \(p\); their product has order
\[
 \operatorname{ord}(\zeta_M\zeta_p)=\operatorname{lcm}(M,p)=Mp .
\]

Write \(F_{\rm rec}\) for the embedded record field and \(F\) for
the cover coefficient field. The pin on the coefficient field
sends the generators of \(F_{\rm rec}\) to their prescribed complex
values in \(F\), so \(F_{\rm rec}\subseteq F\).
For the marked CM lattice \(\Lambda\) the field conditions read
\[
 k\big(j(\Lambda), W(P):P\in (Mp)^{-1}\Lambda/\Lambda\big)\subseteq F .
\]
Here \(k\) denotes the marked imaginary quadratic field and \(W\) its Weber function. 

For each automorphism
\(\sigma\) of \(\C\) fixing \(F\), these values are fixed. By the main theorem of complex multiplication in its ray-class form \cite{cm-reciprocity}, the Weber values at level \(Mp\) lie in its ray field, contained in \(F\) by Lemma~\ref{lem:marked-original-geometric-field}. The equality of \(j\)-invariants gives an isomorphism to the pullback. Since \(F\) contains the CM field, it respects the specified CM action. Let \(P_0\) generate the free rank-one \(\mathcal O_k/(Mp)\)-module of torsion. Write \(f\) for the isomorphism to the pullback and \(W\) for its Weber quotient. Since \(W\) is the quotient by the CM units, the equality of the two readings gives
\[
 W(f(P_0))=W(P_0^\sigma)
 \quad\Longrightarrow\quad f(P_0)=uP_0^\sigma
                  \quad(u\in\mathcal O_k^\times).
\]

The map \(u^{-1}f\) now matches \(aP_0\) for every \(a\in\mathcal O_k/(Mp)\) by CM-linearity. This isomorphism may depend on \(\sigma\). For any \(m\mid Mp\) the
same isomorphism works at level \(m\), since the point with
label \(a/m\) has label \(((Mp/m)a)/(Mp)\) at the original
level.
\label{cor:split-level-sections-1}
\end{proof}

\begin{proposition}[A fixed split refinement for the short theta rules]\label{prop:split-level-sections}
For the split prime and the marked core \(J\) of Lemma~\ref{hyp:split-section-refinement},
one can choose the same cover, coefficient field and weight-zero
determinant character for every number \(n\) of torsion arguments.
For \(n\ge r_j=p_j+q_j\) the Fourier coefficient is obtained from
the basis test below, with the inverse character under a change of
marking. If \(h_n\) denotes the coefficient attached to \(n\) torsion arguments, the rank obstruction gives
\[
 n<r_j\quad\Longrightarrow\quad h_n=0.
\]
\end{proposition}

\begin{proof}
For the component comparison the refined auxiliary level is
one, the rank is positive, and the marked lattice supplies a
rank-one definite datum for the same host. In indefinite
signature the rank is at least three. Under these hypotheses
the CM comparison for the unitary period spaces says that the
pullback isomorphism fixing the torsion labels and the host
root induces the identity permutation of the period
components. Consequently an algebraic constant in the
function field of a component is fixed by every
\(\sigma\in\operatorname{Aut}(\C/F)\) under its generic
embedding. Indeed, if an algebraic constant \(a\) had degree greater than one, separability and extension of embeddings to \(\mathbb C\) would give
\[
 \sigma\in\operatorname{Aut}(\mathbb C/F),\qquad \sigma(a)\ne a,
\]
contrary to the asserted invariance. Thus \(F\) is relatively algebraically closed in
each component function field. The component idempotents and
the prescribed values in \(F\) now give the componentwise
weight-zero algebraic rules.
\label{cor:split-level-sections-2}

\proofstep{The determinant character.}
Here one must put a nontrivial character on that refinement.
Let \(\mathbb A_{E,f}^1\) be the norm-one subgroup of the
finite ideles. Hilbert 90 gives
\[
 1\longrightarrow\mathbb A_{F,f}^\times
 \longrightarrow\mathbb A_{E,f}^\times
 \xrightarrow{\,h:y\mapsto y/\bar y\,}\mathbb A_{E,f}^1
 \longrightarrow1 .
\]
The finite character \(d_p\) read from \(\delta_p\) is trivial on
the left-hand subgroup. In taking the finite part its
archimedean kernel is also trivial, by divisibility of the
connected complex idele factors and finite order of the
character. Hence \(\xi_p(h(y))=d_p(y)\) is independent of the
lift \(y\). A field-valued complex character on this
subgroup extends as a \(\C^\times\)-valued character of the
abelian group, by divisibility of \(\C^\times\). Only its values
on norm-one determinants will be used. A rational norm-one
determinant has a rational lift by Hilbert 90, so \(\xi_p\)
is trivial there.

At principal level the local congruence
version provides an integral lift congruent to one at \(p\).
At the odd split special prime the componentwise involution interchanges factors, so the lift \(z=(u,1)\) gives
\[
 z/\bar z=(u,1)/(1,u)=(u,u^{-1}).
\] The tame
character therefore gives one on the required level kernel,
and its nontrivial residue value at \(p\) survives on the
determinant quotient. Its values on norm-one determinants are
the original finite-order values in the coefficient field.
The function \(g\mapsto\xi_p(\det g)\) is constant on the
corresponding marked components, by the rational and level
invariance. The preceding component comparison gives a
weight-zero rule with these values, whose same positive power
is one on every component.
\label{cor:split-level-sections-3}

\proofstep{The finite Fourier coefficient.}
The finite Fourier expansion must retain the character just obtained on the determinant.

Write \(r_j=p_j+q_j\) and consider \(n\) torsion arguments on
this cover. For \(n\ge r_j\) fix a basis \(y\) of its second
split half module of rank \(r_j\). On an \(n\)-tuple \(x\)
whose first \(r_j\) entries form a basis \(g y\), set
\[
              h_n(x)=\psi(g^\dagger)^{-1}.
\]
Here \(\psi\) is the character on automorphisms of the first
half induced by the nontrivial determinant character, and
\(g^\dagger\) is the adjoint between the two dual halves.
Set \(h_n=0\) elsewhere. Over the finite residue quotient these
modules are free and a spanning \(r_j\)-tuple is a basis,
as can be checked at every local factor by its determinant.
The basis change \(g\) is unique. Thus \(h_n(y,0)=1\). The inverse character under change of marking follows from reversal of composition by adjunction:
\[
 (g'g)^\dagger=g^\dagger(g')^\dagger .
\]

Take the finite Fourier
sum of \(h_n\) against the additive character of the level
pairing in the \(n\) arguments. Its coefficients lie in the
field just specified. This is the finite coefficient of the
short theta rule at this order. When \(n<r_j\) its spanning
condition has empty hypothesis and we take the zero coefficient.
The weight-zero twist section, the underlying cover and its
coefficient field were fixed before \(n\). Only the coefficient
of the theta rule changes. Finite Fourier evaluation and the
harmonic chart transformation compute these short rules on the
same cover at all orders, with the same determinant twist.
\label{cor:split-level-sections-4}
\end{proof}


\begin{lemma}[One level action for both orientations]
Let \(\bar A_g^{(1)},\bar A_g^{(2)}\) be the two induced operators of the finite row group on the core half-level quotient \(\Lambda/\mathfrak m_b\Lambda\) at lift \(j\). With the common level-form identification,
\[
 \bar A_g^{(1)}=\bar A_g^{(2)}\quad\text{on }\Lambda/\mathfrak m_b\Lambda\qquad\text{for every }g.
\]
This action intertwines that on the level forms.
\end{lemma}
\begin{proof}
Write \(\Lambda\) for this lattice and \(A_g\) for the core level
map at \(g\). The congruences of the level comparison give
\[
 \bar A_1=1,\qquad \bar A_g\bar A_h=\bar A_{gh}
       \quad\text{on }\Lambda/\mathfrak m_b\Lambda .
\]
In particular \(\bar A_{g^{-1}}\) is the two-sided inverse
of \(\bar A_g\). This action is defined on \(J\) before either
extension. Two integral lifts of \(\bar A_g\) differ by a map
into \(\mathfrak m_b\Lambda\); on the complementary level
sublattice their difference lands in the full level multiple.
Thus they give the same allowed change of level vector. Pullback
on forms depends only on the quotient map, and the geometric
action in the common frame intertwines it with the row action.
\label{cor:completed-division-9}
\end{proof}

At \(v\) use the extension at \(p(v)\). Its own prime
differs from the residue characteristic, including when
the residue characteristic is \(p_d'\). The place
\(\bar v\) uses the same orientation as \(v\).
Consequently the short-cover laws, the
determinant comparisons and the one level action apply
on the indicated parts of this single numerical point.
No scale cutoff is used in this selection. The linked row
hypotheses and the numerical inequalities for section counts
have their own thresholds.
\label{cor:completed-division-10}


\begin{lemma}[Common form coordinates on the two runs]
At each lift \(j\), write \(I_j\) for the common arithmetic frame identification. The same inverse coordinate map can be used for both orientations:
\[
 z_j^{(1)}=z_j^{(2)}=I_j^{-1}=:z_j.
\] It reads the
same finite action, good projector and depth scalar at \((v,j)\).
The short sections at \(v\) come from its own prime \(p(v)\).
\end{lemma}
\begin{proof}
Use the phase-one base row of \(P\). Its weights are the positive
entries \(a^+_{j,\tau}\) and the reversed negatives of
\(a^-_{j,\tau}\), of signature \((k+j-1,m-j)\) and central
weight \(n_V\). These are the product order and the light and
short weights of the numerical record.

Let \(I_j\) denote the injective comparison of core cusp forms
with the arithmetic frame. The level quotient determines the
pullback \(A_g\) for each element of the mirabolic group.
It suffices to read the representation \(\rho\) at one coefficient
place, because \(I_j A_g=\rho(g)I_j\) in the common frame.
In particular
\[
 I_jA_gA_h=\rho(g)\rho(h)I_j=I_jA_{gh}.
\]
Injectivity gives composition, and gives linearity by the same
argument. The basis comparison makes the scalar extension of
\(I_j\) invertible.

Take \(z_j=I_j^{-1}\) on this extension, as
in Paragraph~\ref{par:selected-form-coordinates}. For the
arithmetic projector \(e_{\rm good}\) and a rational class-sum
map \(T_\mu\) in these coordinates, the corresponding form
readings are \(z_je_{\rm good}I_j\) and \(z_jT_\mu\).
Thus
\[
 I_j(z_je_{\rm good}I_j)=e_{\rm good}I_j,\qquad
 I_j(z_jT_\mu)=T_\mu .
\]
Use this one isomorphism for the ambient lattice and the scaled
and unscaled short sections at \(p(v)\). The section lattice in
the good row is the projector image of the theta lattice at
\(\mu(v)\).

The embedding of the host integers into the point field factors
through the integral coefficient ring of the core, and agrees
with complex conjugation on the host involution. It therefore
reads the same half-level and error generators and the same
scalar in the form and depth comparisons. Its integral exponent
is the depth in the local section estimate. The remaining
parameters \(f,m,r,k,x,\epsilon,b',X''\), the type primes and
the determinant and cutoff data are those of the common numerical
point. Only the local short-section spans depend on the
distinguished prime. \label{cor:completed-division-11}
\end{proof}


\begin{proposition}[One remainder before choosing the cut]
At the common point of Proposition~\ref{prop:completed-division-pair},
assume the linked rows of Proposition~\ref{prop:linked-comparison-rows}
and the section counts of Proposition~\ref{prop:slice-rank-shares}.
Use the simultaneous short basis of
Lemma~\ref{lem:slice-short-basis-before-localization} on the cut by
the selected light columns, with \(C_{n_W}=10\), \(\log g_i\le9L+1\), \(C_P\ge1000\)
and \(C_P^{20}\le frL\). At sufficiently large scale the section
remainder on each lift of either orientation satisfies
\[
                 \langle r\rangle\le R_j fmb\delta_2 .
\]
The scale threshold is common to the orientations and lifts.
At each place the selected remainder precedes the choice of the
cut and its integer depth.
\end{proposition}
\begin{proof}
Use the core coefficient bound \(C_K\) in \(S_\kappa\le C_K kfb\)
in the scale inequalities \eqref{eq:slice-remainder-scales}.
Write \(X_{\rm net}(20,C_K)\) for a scale beyond which the size
inequalities of Proposition~\ref{prop:slice-evaluation-net} hold
with exponent \(20\) and this core bound. Increase the common
scale threshold \(x_1\) to satisfy
\[
 x_1\ge\max\{10^{60},X_{\rm net}(20,C_K)\}.
\]
Read the half-error divisibility on \(J\), and the nonzero short
sections on the selected full extension. The full, short and
light weights and the index depth are those of
Section~\ref{subsec:paired-row-indices}.

At \(p=p(v)\) fix the split marking, once for all three weight spaces. The input level and its special character are
\[
 g_0g_1p,\qquad g\longmapsto\xi_p(\det g),\qquad \xi_p(y/\bar y)=d_p(y).
\] Read them in the core weight spaces at
\(\kappa_j,w_{F,j},w_{L,j}\) by the common form isomorphism.
The incidence chart is the one for the light columns of this
run on the core. Its full level in the light-weight inequalities
is \(g_0g_1\). Select the \(d_s\) scaled short sections before localization by
Lemma~\ref{lem:slice-short-basis-before-localization}, using independent
evaluation on the geometric generic intersection. The shifted
discriminant of this basis is computed in \eqref{eq:slice-discriminant-order}.

If \(A\) is the selected square evaluation minor over the generic
field \(K\), then for each branch field \(K_\zeta/K\) one has
\[
                    0\ne\det A\in K\hookrightarrow K_\zeta .
\]
Thus the same sections give full rank on every branch. Their local
integral index still contributes to \(I_\lambda\). The choice by rank in the
boundary and complementary spaces preserves the light columns
and, before the final lift, the same auxiliary section
\eqref{eq:selected-eta-product}. Hence the finite valuation
test, the two sections for the height calculation and the
determinants concern this one family with the new choice of sections. The evaluation
net gives the second norm on the same rational spaces and
integral section lattices. Changing this norm in the
determinant proof preserves the index on those modules.


Let \(e\) be the branch ramification index. At its reference point the weight-\(w_F\) section \(g_0\) fixes the valuation offset
\[
 c_\zeta=v(g_0(y_\zeta)).
\]
It cancels against the denominator in the quotient reading.
The Gauss value \(\nu_\zeta\) of the generic light row is the
minimum over its light columns, giving the shifted order
\[
              e\bigl(v_{y_\zeta}(F)-T\nu_\zeta\bigr).
\]
This uses the normalized point valuation of the depth estimate.
The vertical, horizontal and archimedean comparisons in
\eqref{eq:slice-height-disc} bound the discriminant and auxiliary
costs on this slice. The remaining two costs are its boundary
determinant and complementary CM determinant heights. Their
ranks give the \(s^2/T\) savings in \eqref{eq:geometric-four-terms}.
At the definite lift the auxiliary cost and complementary rank
are zero. The four bounds give the claimed inequality with the
chosen finite net and the same polynomial bounds at \(frL\)
for every \(j\) and both \(p\).
\label{cor:completed-division-12}
\end{proof}

Extending the form frame to the point field now puts this remainder
estimate in the coordinates of the paired index bound. The
restriction and depth argument uses that bound on the same
lattice and level action. \label{cor:completed-division-13}

\subsection{The joint estimate on the completed family}
\label{subsec:joint-family-estimate}

The two parameter runs now give local inequalities on one realized division pair. Returning to the original period coordinates and summing over places will put them in the same determinant estimate as the cohomological numerator.

At each scale retain the field and the height exponent of
Lemma~\ref{lem:window-parameter-choice}. The original-host estimate in
Proposition~\ref{prop:actual-rank-budget} applies to all the required
twists. Lemma~\ref{lem:strict-budget-after-cut} restricts its sum to
the good types, and Proposition~\ref{prop:weighted-dual-cost} bounds
the localized numerator on the linked row. We now compare its period
lattices on the two covers of that same row.
\begin{definition}[Indices on a pair of conjugate places]\label{def:joint-pair-indices}
For a fixed completed member of Definition~\ref{def:completed-member}, write \(\omega_w=\log\mathrm Nw/[K:\mathbb Q]\) for the coefficient
place weights, \(e_w\) for absolute ramification and
\(\mathcal E=x^\epsilon(B_0+N_c)\).
At a twisted place \(w\mid p\), let \(I_{wj}\) be the sum of the
two signed transport indices of lift \(j\) at \(w,\bar w\). The
error before passing to the restricted lattice \(r^0_{j,w}\), forced term \(F_j(w)\) and integer
least exponent \(J_{j,w}\) are as in Section~\ref{sec:word-division}.
\end{definition}

\begin{lemma}\label{lem:joint-family}
Fix the family with the parameter choice of
Lemma~\ref{lem:window-parameter-choice} and the linked base
and refined rows of Proposition~\ref{prop:linked-comparison-rows}.
Use the good types of Lemma~\ref{lem:strict-budget-after-cut} on the
original host and the integral numerator of
Proposition~\ref{prop:weighted-dual-cost} on these rows. There are \(C_0\ge0\) and \(X_*\) such that, at
every scale \(x\ge X_*\) in the family's domain, the base row of
each compatible completed member of Definition~\ref{def:completed-member}
satisfies, for \(C\ge C_0\),
\begin{align}
 \sum_j I_{wj}
 &\le D_0\left(\sum_j(2e_wF_j(w)+e_w r^0_{j,w})
                         +2e_w/p^2\right),\label{eq:joint-parent-local}\\
 r^0_{j,w}&\ge0,\qquad
 \sum_{w\ {\rm twisted}}\omega_w(e_w r^0_{j,w}+J_{j,w})
                \le C\mathcal E.\label{eq:joint-parent-global}
\end{align}
On the final members of Construction~\ref{hyp:family-prescribed-realization}, use the reduced chart and product lattices of Lemmas~\ref{hyp:slice-chart-branch-lattices} and~\ref{hyp:slice-product-lattices} on the marked rows. The transport of Lemma~\ref{lem:parent-transport-cost} then compares the lattices at each compatible completion.
The constants are uniform in the translates, words, lifts and
places of the family.
\end{lemma}

\begin{proof}
\proofstep{Choose a comparison that reads the parent row.}
Choose the comparison in the completed division pair with
the prescribed identification of coefficient embeddings, series
reading and geometric core, as in Paragraph~\ref{par:linked-family-rows}.
The primitive and divisor identifications hold on each such
summand. Fix a compatible completed member \(P\), and write \(D(P)\) for
its period maps, integral lattices and numerical sizes. The covolume
comparison and the evaluation-grid comparison may choose different
realizations over \(P\). Denote those identifications by \(\iota_{\rm cov}\) and
\(\iota_{\rm ev}\), on the covolume and evaluation realizations
over \(P\), respectively. Their period-and-lattice records satisfy
\[
         \iota_{\rm cov}(D_{\rm cov})=D(P)
                                =\iota_{\rm ev}(D_{\rm ev}).
\]


More explicitly, let \(C_{\rm cov}(P)\) and \(C_{\rm ev}(P)\) denote the choices of auxiliary comparisons inducing the fixed data \(D(P)\). The constructions on the fixed row give an element of each set. Any pair in
\[
                  C_{\rm cov}(P)\times C_{\rm ev}(P)
\]
gives a simultaneous choice over \(P\). Apply the two inequalities there. Their
row, weight, divisor and coefficient variables then coincide. For each further restriction
with threshold \(X_i\), work on \(x\ge\max_i X_i\) in the
domain of the chosen family.
The rank, divisor data and choice of orientations in
Section~\ref{subsec:completed-prime-runs} now apply there.
\label{cor:joint-family-estimate-1}

\proofstep{Descend the class sums to one coefficient field.}
Let \(A(y)\) denote the complex form of a vector in the lift's
coefficient frame over \(K\), and let \(\iota:K\hookrightarrow
\mathbb C\) be its embedding. The images of a frame basis are
complex linearly independent by the form isomorphism and its
injective reading on the generating lowest vector. We need the
class-sum map at the actual tuple level over this field. Its
rationality uses the ray fields from Paragraph~\ref{par:linked-family-rows}.
Write \(K_{2}\) for the imaginary quadratic subfield. Since it has degree two and contains the indicated square root,
\[
 \mathbb Q(\sqrt{-d})\subseteq K_{2},\qquad [\mathbb Q(\sqrt{-d}):\mathbb Q]=[K_{2}:\mathbb Q]=2,
\]
and the two fields coincide.

At tuple level \(M\ge3\), the Weber values of
the level torsion lie in the embedded ray field by complex
multiplication. With the root of \(-d\) fixed, the Weber comparison
gives the CM curve with the required level structure. The level
of the cover divides \(M\), since its auxiliary level is one.
The component comparison using the rank-one CM datum then gives
geometric irreducibility of its components over the coefficient
field, equivalently relative algebraic closure of that field in
their function fields. Use this comparison when choosing the half
levels and their covers on the normalized reference frame.

On
each tuple take the finite class representatives and their
stabilizer masses in the new level. In an integral frame its local factor takes the form
\[
         K_{t,v}(z)=\mathbf1_{A_{t,v}}(z)
                         \prod_\nu \chi_{\nu,t,v}(g_\nu(z))^{a_\nu},
\]
where \(A_{t,v}\) is the selected finite support and the exponents
are integral. The admitted character values lie in \(\iota K\).
Thus each local factor takes its values there, including zero
off the support. The finite chart products therefore have the
same coefficient field at all tuples. The original one-map comparison also extends from
the reference frame to these tuple levels on the prescribed
quotient primes. The change to the split character reading
preserves its levels, individual chart maps and extremal vector.
\label{cor:joint-family-estimate-2}

Fix a lift of polynomial degree \(d_W\). With the same class factor \(c_\alpha\) in both presentations, their ordered pair of factors is
\[
 (u_1c_\alpha,u_2c_\alpha)=c_\alpha(u_1,u_2).
\]
To see the uniformity in this assertion, compare both readings
of a single class in a standard frame. Identifying the adelic
kernel with the short theta series gives an orientation
\(\delta\in\{1,-1\}\). Evaluation of its Fock coefficient on
the chosen extremal vector gives a phase \(a\), and comparison
on the cover gives a phase \(c_0'\). A sign change of the
polynomial argument acts by its total degree. Consequently the
multiplier for a presentation is
\[
                 u=a(-\delta)^{d_W}/c_0' .
\]

The positive Lie multiplier is
\(L=(\sqrt{2\pi\xi_{\mathbb R}})^{-d_W}\), with
\(\xi_{\mathbb R}>0\) in the prescribed vacuum metric.
The remaining factor \(c_\alpha\) compares the normalized
frames on a component. Compute it once on the reference pair
using its rational translation law for the density. Change of
a marking in the same component gives the same value by the
chart transformation rule. The orientation, extremal-vector
phase and cover phase above are fixed before the tuple and
its class vector. Summing the single-class identity with the
stabilizer masses thus gives the same \(u\) for every tuple
and \(\lvert u c_\alpha\rvert=1\).
\label{cor:joint-family-estimate-3}

For a fixed finite argument the two readings of the series
on a source vector give
\[
 u_1c_\alpha L m_g a_g=u_2c_\alpha L m_g a'_g .
\]
The Lie factor \(L\) is positive. Cancelling \(u_1c_\alpha L\)
gives \(m_g a_g=(u_2/u_1)m_g a'_g\), also when the component
density \(m_g\) vanishes. For the second chart, the determinant character on the
cover and the finite twist have finite orders. Their product
of orders clears both multipliers. With the coefficient-field
character values on the half level, the class sum then
satisfies the descent law on the cover. Its weight is dominant
with positive last weight, so the descended section is cuspidal.

By scalar extension for cusp sections there are finitely many base cusp sections \(f_i\) and coefficients \(a_i\) such that the descended section \(f\) satisfies
\[
 f=\sum_i a_i f_i,\qquad a_i\in\iota K .
\] Read each
base section in the form frame and lift its coefficient
through \(\iota\). Translating that section by the level action
gives the finite translate of its adelic reading with the
prescribed character factor. Incorporating the tuple scalar in its
coordinates gives, for every tuple \(t\) and source \(a\),
\[
       \mathrm{ph}\,\Theta_t(a)=\rho A(y_{t,a}),\qquad
       \rho=u_2/u_1,\quad y_{t,a}\in K^{n_W}.
\]
The equality on the lowest-vector evaluations gives the equality
of forms by injectivity on that space. This uses the same
finite group action for the two readings.
\label{cor:joint-family-estimate-4}

To descend the common scalar, use the good projector \(e\) at
the reference tuple. Its rational lift matrix is bijective by
the base comparison, and the target has dimension \(D_0>0\).
Lift the vector with all coordinates one by its inverse. The
resulting good theta reading is nonzero by injectivity of the
target coordinates.

At this tuple the comparison gives rational
coordinates for the phased good lift. Applying the good projector
to the preceding equality therefore gives
\[
           A(y_1)=\rho A(e y)\ne0.
\]
Complex linear independence gives
\(\rho=\iota((y_1)_h/(e y)_h)\) at any coordinate with
\((e y)_h\ne0\). Absorb this field element into each
\(y_{t,a}\). The phase in the selected comparison is one.
For each tuple restrict the resulting rational map to its level
source. Its preimages under \(A\) are unique, hence additive
and \(K\)-linear by semilinearity of both readings. This
produces the rational class-sum operator \(\theta\) at every
assigned pair \((w,j)\), even for a tuple with zero class sum.

\proofstep{Verify equivariance and restriction on the good space.}
Write the central good idempotent as \(\sum_g a_g[g]\), with
actions \(e_X,e\) on the source and target. The descended
class-sum map intertwines each \(g\), so
\[
     \theta e_Xx=\sum_g a_g g\,\theta(x)=e\theta(x).
\]
Centrality gives the same identity after restriction to the level
character. To project in the source level metric, take a basis
\(k_1,\ldots,k_h\) of \(\ker\theta\) and its Gram matrix \(G\),
using the convention that the Hermitian form is linear in its
second argument. The restricted matrix is invertible at the
selected embedding, since the metric is positive definite there.
Thus the correction has coordinates in the field:
\[
        {\rm pr}(x)=x-\sum_{i=1}^h a_i k_i,\qquad
        Ga=(\langle k_j,x\rangle)_j .
\]

The correction satisfies \(\theta({\rm pr}(x))=\theta(x)\). The base lift is surjective
onto the good rational space. For lattice membership, use
Section~\ref{subsec:paired-row-indices} on this same rational
map. On a tuple generator supported on its class the
inverse-character trace over its stabilizer belongs to the
level lattice. Its class-sum image is the chosen generator
image times an integral character unit. The span comparison gives equality on generic fibres, and at each of the indicated places it gives the integral inclusion
\[
 \mathcal X_{3,w}\subseteq(e\theta(M_t))_w .
\] The two places \(w,\bar w\) have
the same tuple and residue characteristic.
\label{cor:joint-family-estimate-5}

\proofstep{Apply the geometric and arithmetic estimates to one record.}
The local transport inequality and the weighted cohomological numerator have now been obtained. Their sum can be used in the determinant only after checking that their ranks, place labels and lattices all refer to the chosen member.

Use the simultaneous choice of the outer, middle and last
parameters in Section~\ref{subsec:window-points} for the
finitely many family conditions. For the resulting family
and \(C_H\), the run and rationality arguments have cutoffs
\(x_R,x_Q\). The geometric, rank, twist, character-share,
bad-type-share, exponent and numerical size estimates each
have a lower bound on their main comparison constant.
Choose \(c_0\) at least the maximum of these seven bounds
and \(20\), and put
\[
 c=(c_0,14,66c_0+4),\qquad C_b=14+(66c_0+4)+4,\qquad
 C_E=4c_0^2+25c_0+C_b+44 .
\]
The last two entries of \(c\) meet the two bad-prime allowances.
Write \(x_i(c)\), \(1\le i\le7\), for the thresholds of these estimates
at the chosen \(c\), and let \(x_{\rm fam}\) dominate the preceding family
cutoffs.

Take
\[
 X_*=\max\{x_{\rm fam},x_R,x_Q,3c_0,x_1(c),\ldots,x_7(c)\}.
\]
At any point beyond this threshold choose the two runs on its common
geometric core. To build the record \(O\) of
Section~\ref{subsec:completed-prime-runs}, use the second run at the
first distinguished characteristic and the first run everywhere else.
Its section maps and lattices are pulled back through the same
cusp form frame. Its depth uses the canonical half ideals, the
actual error generator and the point's weight data. Every
one of the seven estimates applies to every pair of runs and
record satisfying these identifications. We apply each to
these particular runs and this \(O\).
\label{cor:joint-family-estimate-6}

Here are the hypotheses of those applications.
\begin{enumerate}
\item The numerical sizes come from the parameters of this family
at \(x\). After increasing \(X_*\), the numerical conditions used for this point are simultaneous:
\[
 x\ge3c_0,\qquad C_{n_W}=10,\qquad \mathcal E\ge0 .
\]
The two primes and the level in this application are those of the point. The level and
cutoff comparisons of Section~\ref{subsec:completed-prime-runs}
give the bounds on the level squared and the cutoff \(Y_0\)
in the prime sum.
\item Use the geometric section theorem of that section with
the orientation prescribed at each place. Its short sections
span, their scaled lattice lies in the ambient lattice, and
its remainder is chosen before restriction to a subspace.
Pulling its estimate to the selected good subspace gives the
index estimate of Section~\ref{subsec:paired-row-indices}.
Use its four depth bounds at the pair \((C_E,C_H)\) for the
ordinary, ramified, dyadic and level-prime cases. If \(a_w\) is the scaling scalar at the place under consideration, the forced content used in this application is read at the same place and with the same normalization:
\[
 \text{forced content at }w=v_w(a_w).
\]
\item For the rank comparison take the integral lattice \(J\)
of base-character generators of the selected run. The class representatives of
the base theta kernel give the spanning list. Indeed evaluation of a generator on its own finite
representative recovers the class series, and tracing that
representative over its stabilizer supplies every class in
the list. The deletion-character values are in the point's
field. The off-prime comparison for the used run identifies the
localized short theta generators inside \(J\). The latter is
finitely generated and lies in the level character space.
The finite good projector commutes with the action there.

For the character \(\chi\) and a vector \(v\) in its
space, the unnormalized trace and its projected image are
\[
      \sum_{u\in U}\chi(u)^{-1}uv=|U|v,\qquad
      e\sum_{u\in U}\chi(u)^{-1}uv=|U|ev .
\]
The short sections span the character space over the
fraction field. Since \(|U|\ne0\), their projected
traces give the full generic good space of \(eJ\). Apply this also at the conjugate place, using
the same off-prime choice. \label{cor:joint-family-estimate-7}
\item Now use the simultaneous test for the chosen twisting
character with this full-rank baseline. At a place \(v\) of the chosen pair, Section~\ref{subsec:paired-row-indices} bounds the excess over the full-rank baseline:
\[
 [M_g:\mathcal X_3]_v-[M_g:eJ]_v\le e_v\epsilon_{vj}.
\]
On a pair the sum of both excesses over all lifts is at most
\(2/p^2\). Its application uses the integral inclusion above
and the prescribed character at both places.
\item For the character share the common action on the cusp
forms identifies the trace comparison with that for the
actual weight. For the discarded share the two central
finite-type idempotents act in the same representation.
Equation~\eqref{eq:row-bad-share} bounds the discarded rank by
the regular-representation trace times the multiplicity bound.
Together the two estimates give
\[
 (1-y)R_j\le|U|d_0,\qquad (1-y)|U|\le I,\qquad
 (1-\delta)d_0\le D_0,\qquad y\le1/8,\quad\delta\le1/4 .
\]
Here \(I\) is the level index. Hence \(R_j\le3 I D_0\).
\item Finally the integral exponent comparison at \(p\) uses
the actual pair of lift lattices in the record. The inclusion
by the rational-prime power \(p^{B_j(p)}\) gives
\(0\le J_{j,w}\le e_w B_j(p)\). Grouping the places above each rational prime and using \eqref{eq:division-least-sum} therefore gives
\[
 \sum_{w\mid p}e_w\omega_w=\log p,\qquad \sum_{w\ {\rm twisted}}\omega_w J_{j,w}\le32\mathcal E .
\] The depth and exponent
arguments use the same coefficient field as the signed lengths.
\end{enumerate}

The restriction, signed length and conjugation calculations in
Section~\ref{subsec:paired-row-indices} now apply to \(O\) with
the two class-sum identities and the base inclusion. The character
and central projections in that calculation are the projections
in its pulled-back level action. The lengths and rank therefore
belong to the base row of this very member. By Section~\ref{subsec:parent-row-errors}, at every lift the three contributions to the weighted auxiliary error together obey
\[
 \text{weighted auxiliary error at }j\le C_E\mathcal E\qquad(1\le j\le m).
\] On forgetting the split phase-one comparison,
the normalized row has phase one. Substitution in the period
and numerical readings identifies the matrices, both
lattice indices, \(D_0\), the forced weights and the normalized
exponents \(J_{j,w}/e_w\) with their parent values.
They have these values for every choice of the numerical
allowance \(C\).

Write \(\mathfrak e_j(w)\ge0\) for the auxiliary
transport error of the comparison before descent. The local
bound gives
\[
   B_w:=D_0^{-1}\sum_h I_{wh}
         -2\sum_h e_wF_h(w)-2e_w/p^2
       \le \sum_h e_w\mathfrak e_h(w).
\]
The summands and the right-hand side are read at the selected
completed member. Therefore its defining error
\(e_w r^0_{j,w}=[B_w]_+/m\) is nonnegative and
\(B_w\le\sum_j e_w r^0_{j,w}\), proving
\eqref{eq:joint-parent-local}. For the weighted bound, first
discard the nonnegative least exponents in each auxiliary
transport budget. This gives
\[
 \sum_{w\ {\rm twisted}}\omega_w e_w r^0_{j,w}
 \le \frac1m\sum_h\sum_{w\ {\rm twisted}}
                    \omega_w e_w\mathfrak e_h(w)
 \le C_E\mathcal E.
\]
For the fixed lift \(j\), discard instead the transport term
to obtain \(\sum_{w\ {\rm twisted}}\omega_w J_{j,w}\le C_E\mathcal E\).
Their sum proves \eqref{eq:joint-parent-global}, with
\(C_0=2\max(C_E,0)\), at every member above the common threshold. \label{cor:joint-family-estimate-8}
\end{proof}

